\documentclass[
10pt, % Main document font size
a4paper, % Paper type, use 'letterpaper' for US Letter paper
oneside, % One page layout (no page indentation)
%twoside, % Two page layout (page indentation for binding and different headers)
headinclude,footinclude, % Extra spacing for the header and footer
BCOR5mm, % Binding correction
]{scrartcl}
%%%%%%%%%%%%%%%%%%%%%%%%%%%%%%%%%%%%%%%%%
% Arsclassica Article
% Structure Specification File
%
% This file has been downloaded from:
% http://www.LaTeXTemplates.com
%
% Original author:
% Lorenzo Pantieri (http://www.lorenzopantieri.net) with extensive modifications by:
% Vel (vel@latextemplates.com)
%
% License:
% CC BY-NC-SA 3.0 (http://creativecommons.org/licenses/by-nc-sa/3.0/)
%
%%%%%%%%%%%%%%%%%%%%%%%%%%%%%%%%%%%%%%%%%

%----------------------------------------------------------------------------------------
%	REQUIRED PACKAGES
%----------------------------------------------------------------------------------------

\usepackage[
nochapters, % Turn off chapters since this is an article        
beramono, % Use the Bera Mono font for monospaced text (\texttt)
eulermath,% Use the Euler font for mathematics
pdfspacing, % Makes use of pdftex’ letter spacing capabilities via the microtype package
dottedtoc % Dotted lines leading to the page numbers in the table of contents
]{classicthesis} % The layout is based on the Classic Thesis style

\usepackage{arsclassica} % Modifies the Classic Thesis package
\usepackage{ulem}
\usepackage{graphicx}
\usepackage{float}
\usepackage{algorithm2e}
\usepackage[export]{adjustbox}
\usepackage[utf8]{inputenc}  
\usepackage[english,frenchb]{babel} 
\usepackage[T1]{fontenc} % Required for accented characters
\usepackage{multicol}
\usepackage{graphicx} % Required for including images
\graphicspath{{Figures/}} % Set the default folder for images

\usepackage{enumitem} % Required for manipulating the whitespace between and within lists

\usepackage{lipsum} % Used for inserting dummy 'Lorem ipsum' text into the template


\expandafter\def\csname ver@subfig.sty\endcsname{}


\usepackage{svg}
\usepackage{hyperref}

\usepackage{amsmath,amssymb,amsthm} % For including math equations, theorems, symbols, etc

\usepackage{varioref} % More descriptive referencing

%----------------------------------------------------------------------------------------
%	THEOREM STYLES
%---------------------------------------------------------------------------------------

\theoremstyle{definition} % Define theorem styles here based on the definition style (used for definitions and examples)
\newtheorem{definition}{Definition}

\theoremstyle{plain} % Define theorem styles here based on the plain style (used for theorems, lemmas, propositions)
\newtheorem{theorem}{Theorème}
\newtheorem{lemma}[theorem]{Lemme}
\newtheorem{proposition}[theorem]{Proposition}

\theoremstyle{remark} % Define theorem styles here based on the remark style (used for remarks and notes)

%----------------------------------------------------------------------------------------
%	HYPERLINKS
%---------------------------------------------------------------------------------------

\hypersetup{
%draft, % Uncomment to remove all links (useful for printing in black and white)
colorlinks=true, breaklinks=true, bookmarks=true,bookmarksnumbered,
urlcolor=webbrown, linkcolor=RoyalBlue, citecolor=webgreen, % Link colors
pdftitle={}, % PDF title
pdfauthor={\textcopyright}, % PDF Author
pdfsubject={}, % PDF Subject
pdfkeywords={}, % PDF Keywords
pdfcreator={pdfLaTeX}, % PDF Creator
pdfproducer={LaTeX with hyperref and ClassicThesis} % PDF producer
}

\hyphenation{Fortran hy-phen-ation} 

\title{\normalfont{FTML Exercices 1 solutions}\\
\large{Pour le 23 avril 2026}
} 
\author{\spacedlowsmallcaps{}
} 
\date{}
\begin{document}
\renewcommand{\sectionmark}[1]{\markright{\spacedlowsmallcaps{#1}}}
\lehead{\mbox{\llap{\small\thepage\kern1em\color{halfgray} \vline}\color{halfgray}\hspace{0.5em}\rightmark\hfil}} 
\pagestyle{scrheadings}
\maketitle 
\setcounter{tocdepth}{3} 

\tableofcontents

\section{\large\color{Blue}Probabilités et statistiques}

\subsection{\color{MidnightBlue}P1}

\subsubsection{Enoncé}
\label{par:enoncé}

Calculer l'espérance et la variance des variables aléatoires réelles suivantes.
\begin{itemize}
    \item $X_1$ de loi uniforme sur $[0, 1]$.
    \item $X_2$ de loi uniforme sur $[-1, 1]$.
\end{itemize}

\subsubsection{Solution}


$X_1$ et $X_2$ sont des variables à densité. Pour tout $x$ dans leur support,
$p_{X_1}(x) = 1$, et $p_{X_2}(x) = \frac{1}{2} $.

\begin{equation}
    \begin{aligned}
        \label{eq:}
        E[X_1] &= \int^{1}_{0} x p_{X_1}(x) dx \\
        &= \int^{1}_{0} x  dx \\
        &= [ \frac{x^2}{2} ]^1_0\\
        &= \frac{1}{2} 
    \end{aligned}
\end{equation}

\begin{equation}
    \begin{aligned}
        \label{eq:}
        E[X_2] &= \int^{1}_{-1} x p_{X_2}(x) dx \\
        &= \int^{1}_{-1} x \frac{1}{2}  dx \\
        &= \frac{1}{2}[ \frac{x^2}{2} ]^1_{-1}\\
        &= 0
    \end{aligned}
\end{equation}

On dit que $X_2$ est \textbf{{centrée.}} 

\subsection{\color{MidnightBlue}P2}

\subsubsection{Enoncé}
\label{par:enoncé}

Calculer l'espérance et la matrice de variance-covariance du vecteur aléatoire suivant

\begin{equation}
    Y = (Y_1, Y_2)
\end{equation}

où
\begin{itemize}
    \item $Y_1$ suit une loi de Bernoulli de paramètre $p$
    \item $Y_2$ suit une loi normale $ \mathcal{N}(\mu, \sigma)$.
    \item On suppose $Y_1$ et $Y_2$ indépendantes.
\end{itemize}

\subsubsection{Solution}

Puisque $Y_1$ et $Y_2$ sont indépendantes, on sait que les termes non-diagonaux
sont nuls. Il suffit donc de calculer les variances de $Y_1$ et $Y_2$.
\\

On sait que la variance de $Y_2$ est $\sigma^2$ car c'est une loi normale (
\url{https://fr.wikipedia.org/wiki/Loi_normale}).
\\

Calculons la variance de $Y_1$ (\url{https://fr.wikipedia.org/wiki/Loi_de_Bernoulli} ). Son espérance vaut $E[Y_1] = 1\times p+0\times (1-p) = p$.  Pour calculer la variance on utilise le théorème de transfert.

\url{https://en.wikipedia.org/wiki/Law_of_the_unconscious_statistician} 

Cependant, il ne faut pas oublier qu'il s'agit d'un théorème et non d'une définition.

Si $X$ est une variable réelle à densité $p_X$, alors sous certaines hypothèses (qui sont vérifiées ici)
\begin{equation}
    E[f(X)] = \int_{x\in \mathbb{R} } f(x) p_X(x) dx 
\end{equation}

Appliqué à $Y_1$, cela donne:
\begin{equation}
    \begin{aligned}
        \label{eq:}
        Var[Y_1] &= E[(Y_1-E[Y_1])^2]\\
        &= (1-p)^2\times p + (0-p)^2\times (1-p)\\
        &= p(1-p)(p+1-p)\\
        &= p(1-p)
    \end{aligned}
\end{equation}

Finalement

\begin{equation*}
    Var(Y) = 
\begin{pmatrix}
    p(1-p)  & 0 \\
    0  & \sigma^2 \\
\end{pmatrix}
\end{equation*}


\subsection{\color{MidnightBlue}P3}

\subsubsection{Enoncé}
\label{par:enoncé}

Calculer l'espérance et la matrice de variance-covariance du vecteur aléatoire suivant

\begin{equation}
    Z = (Z_1, Z_2)
\end{equation}

où
\begin{itemize}
    \item $Z_1$ suit une loi uniforme sur $[1, 2]$
    \item $Z_2=Z_1^2$.
\end{itemize}

\subsubsection{Solution}


On a $E[Z_1] = \frac{3}{2} $. On utilise le théorème de transfert pour $E[Z_2]$.

\begin{equation}
    \begin{aligned}
        \label{eq:}
	E[Z_2] &= \int^{2}_{1} x^2 p_{Z_1}(x)dx \\
	&= [\frac{x^3}{3} ]^2_1\\
	&= \frac{7}{3} 
    \end{aligned}
\end{equation}

\begin{equation}
    \begin{aligned}
        \label{eq:}
	Cov(Z_1, Z_2) &= \int^{2}_{1} (x- \frac{3}{2} )(x^2- \frac{7}{3} )  dx \\
	&= \int^{2}_{1} x^3- \frac{7}{3} x - \frac{3}{2} x^2  + \frac{7}{2}  dx\\
	&= [ \frac{x^4}{4} ]^2_1 - \frac{7}{3} [ \frac{x^2}{2} ]^2_1 - \frac{3}{2} [ \frac{x^3}{3} ]^2_1 + \frac{7}{2} \\
	&= \frac{15}{4} - \frac{7}{3}\times \frac{3}{2} - \frac{3}{2} \times\frac{7}{3} + \frac{7}{2} \\
	&= \frac{15}{4} - 7 + \frac{7}{2} \\
	&= \frac{15}{4} - \frac{14}{4} \\
	&= \frac{1}{4} 
    \end{aligned}
\end{equation}

\begin{equation}
    \begin{aligned}
        \label{eq:}
	Var(Z_1) &= \int^{2}_{1} (x- \frac{3}{2} )^2  dx \\
	&= \int^{ \frac{1}{2} }_{ -\frac{1}{2} } u^2du\\
	&= [ \frac{u^3}{3} ]^{ \frac{1}{2} }_{- \frac{1}{2}} \\
	&= \frac{1}{12} 
    \end{aligned}
\end{equation}

\begin{equation}
    \begin{aligned}
        \label{eq:}
	Var(Z_2) &= \int^{2}_{1} ( x^2 - \frac{7}{3} )^2  dx \\
	&= \int^{2}_{1} x^4- \frac{14}{3} x^2+ \frac{49}{9}  dx \\
	&= [ \frac{x^5}{5} ]^2_1 - \frac{14}{3} [ \frac{x^3}{3} ]^2_1 + \frac{49}{9} \\
	&= \frac{31}{5} - \frac{14}{3} \frac{7}{3} + \frac{49}{9} \\
	&= \frac{31}{5} - \frac{98}{9} + \frac{49}{9} \\
	&= \frac{31}{5} - \frac{49}{9} \\
	&= \frac{31\times 9 - 49\times 5}{45} \\
	&= \frac{34}{45} 
    \end{aligned}
\end{equation}

As the variance matrix is symmetric, we have:
\begin{equation*}
    Var(Z) = 
\begin{pmatrix}
    \frac{1}{12}   & \frac{1}{4}  \\
    \frac{1}{4}   & \frac{34}{45}  \\
\end{pmatrix}
\end{equation*}

\section{\large\color{Blue}Expected value as a minimization}

\subsubsection{Enoncé}%
\label{ssub:Enoncé}

Soit $X$ une variable aléatoire réelle ayant un moment d'ordre $2$. Montrer que
son espérance $E(X)$ est la quantité minimisant la fonction de variable réelle $t\mapsto E\big( (X-t)^2 \big)$

\subsubsection{Solution}%
\label{ssub:Solution}

All expected values are over $X$. We remark that

\begin{equation*}
    \begin{aligned}
	E\Big[(X-t)^2 \Big]&= E\Big[\big(X-E(X)+E(X)-t\big)^2 \Big]\\
	&=E\Big[\big(X-E(X)\big)^2+2\big(X-E(X)\big)\big(E(X)-t\big)+\big(E(X)-t\big)^2 \Big]\\
    \end{aligned}
\end{equation*}

By linearity, the expected value is separated in $3$ terms.

\begin{itemize}
    \item $E\Big[\big(X-E(X)\big)^2\Big]$
    \item $E\Big[2\big(X-E(X)\big)\big(E(X)-t\big)\Big]$
    \item $E\Big[\big(E(X)-t\big)^2\Big]$
\end{itemize}

We note that the first term $E\Big[(X-E(X))^2\Big]$ does not depend on $t$.

Also,  $(E(X)-t)^2$, is a fixed scalar, and not a random variable, hence:

\begin{equation*}
    E\Big[\big(E(X)-t\big)^2\Big]=\big(E(X)-t\big)^2
\end{equation*}

We also have that
\begin{equation*}
    E\Big[2\big(X-E(X)\big)\big(E(X)-t\big)\Big]=2\big(E(X)-t\big)E\Big[\big(X-E(X)\big)\Big]=0
\end{equation*}

As a consequence, the value that minimizes $E\big[(X-t)^2 \big]$ is $t=E(X)$.
\\

On peut aussi prouver le résultat en développant $(X-t)^2=X^2-2tX+t^2$.

\begin{equation}
\begin{aligned}
    \label{eq:}
    E[(X-t)^2]  &= E[(X^2-2tX+t^2)]\\
		&= E[X^2]-2tE[X] + E[t^2]\\
		&= E[X^2]-2tE[X] + t^2\\
\end{aligned}
\end{equation}

La dérivée par rapport à $t$ vaut $2t - 2E[X]$, qui est négative pour $t\leq E[X]$ et positive pour $t\geq E[X]$, ce qui prouve le résultat.



\section{\large\color{Blue}Calcul différentiel}

\subsection{\color{MidnightBlue}C1}

\subsubsection{Enoncé}
\label{par:enoncé}

Calculer le gradient en tout point des applications suivantes.

$$
f_1 = \left\{
    \begin{array}{ll}
        \mathbb{R}^2 \rightarrow \mathbb{R}  \\
	(x,y)\mapsto 3
    \end{array}
\right.
$$

$$
f_2 = \left\{
    \begin{array}{ll}
        \mathbb{R}^2 \rightarrow \mathbb{R}  \\
	(x,y)\mapsto x^3+\sin y
    \end{array}
\right.
$$

$$
f_3 = \left\{
    \begin{array}{ll}
        \mathbb{R}^2 \rightarrow \mathbb{R}  \\
	(x,y)\mapsto x^3\sin y
    \end{array}
\right.
$$

$$
f_4 = \left\{
    \begin{array}{ll}
        \mathbb{R}^d \rightarrow \mathbb{R}  \\
	x\mapsto ||x||_2^2
    \end{array}
\right.
$$

où $||.||_2$ est la norme euclidienne dans $\mathbb{R}^d$.

\subsubsection{Solution}

$\forall (x,y)\in \mathbb{R}^2$, 
\begin{equation*}
    \nabla_{(x,y)}f_1(x,y)=
\begin{pmatrix}
    0\\
    0
\end{pmatrix}
\end{equation*}

$\forall (x,y)\in \mathbb{R}^2$, 
\begin{equation*}
    \nabla_{(x,y)}f_2(x,y)=
\begin{pmatrix}
    3x^2\\
    \cos y
\end{pmatrix}
\end{equation*}

f4) Soit $x\in \mathbb{R}^d$ et $h\in \mathbb{R}^d$.

Ici on utilise la définition du gradient par la différentielle et le développement limité, mais on arriverait au même résultat avec les dérivées partielles (s'entraîner à le faire).

\begin{equation}
\begin{aligned}
    \label{eq:}
    ||x+h||^2 &= \langle x+h, x+h \rangle\\
    &= \langle x, x \rangle + 2 \langle x, h \rangle + \langle h, h \rangle\\
    &= ||x||^2 + \langle 2x, h \rangle + ||h||^2\\
    &= ||x||^2 + \langle 2x, h \rangle + o(h)
\end{aligned}    
\end{equation}

Ainsi 

\begin{equation}
    \nabla_xf(x) = 2x
\end{equation}

\end{document}
